%% sandbox.tex — Unified test suite for the citemsp package (biblatex backend)
%%
%% Exercises every feature and edge case in a single document:
%%   §1  Plain citations and range compression
%%   §2  Default locators (section / paragraph)
%%   §3  Built-in prefix locators
%%   §4  Mixed locator types in a single call
%%   §5  Range compression with mixed entries
%%   §6  Italic, slanted, and bold-italic contexts
%%   §7  Font size contexts
%%   §8  Scaling (\citemspscale)
%%   §9  Small caps and \emph toggle
%%   §10 Math mode adjacency
%%   §11 Theorem and list environments
%%   §12 Table cells
%%   §13 Long locator strings and many-key calls
%%   §14 Double-slash / extra-slash edge cases
%%   §15 Footnotes
%%   §16 Drop-in \cite replacement and realistic text
%%
%% Build:  latexmk -pdf sandbox.tex
\documentclass[11pt, a4paper]{article}

\usepackage[T1]{fontenc}
\usepackage{lmodern}
\usepackage[margin=2.5cm]{geometry}
\usepackage[dvipsnames]{xcolor}
\definecolor{linkblue}{rgb}{0.0, 0.2, 0.6}

%%% BIBLATEX CONFIGURATION
%%% Default: numeric-comp + biber (full compression).
%%% Uncomment one of the alternatives to test other combinations.
%% A) numeric-comp + biber (default — full compression)
\usepackage[style=numeric-comp, backend=biber, sorting=none]{biblatex}
%% B) numeric-comp + bibtex (compression, lighter backend)
%\usepackage[style=numeric-comp, backend=bibtex, sorting=none]{biblatex}
%% C) numeric + biber (no compression)
%\usepackage[style=numeric, backend=biber]{biblatex}
%% D) numeric + bibtex (no compression, lighter backend)
%\usepackage[style=numeric, backend=bibtex]{biblatex}

\begin{filecontents}{sandbox-refs.bib}
@article{einstein1915,
  author  = {Einstein, Albert},
  title   = {Die Feldgleichungen der Gravitation},
  journal = {Sitzungsberichte der Preussischen Akademie der Wissenschaften},
  year    = {1915},
  pages   = {844--847}
}
@article{schwarzschild1916,
  author  = {Schwarzschild, Karl},
  title   = {\"Uber das Gravitationsfeld eines Massenpunktes},
  journal = {Sitzungsberichte},
  year    = {1916}
}
@book{wald1984,
  author    = {Wald, Robert M.},
  title     = {General Relativity},
  publisher = {University of Chicago Press},
  year      = {1984}
}
@article{hawking1974,
  author  = {Hawking, Stephen W.},
  title   = {Black hole explosions?},
  journal = {Nature},
  volume  = {248},
  pages   = {30--31},
  year    = {1974}
}
@book{misner1973,
  author    = {Misner, Charles W. and Thorne, Kip S. and Wheeler, John A.},
  title     = {Gravitation},
  publisher = {W. H. Freeman},
  year      = {1973}
}
@article{penrose1965,
  author  = {Penrose, Roger},
  title   = {Gravitational Collapse and Space-Time Singularities},
  journal = {Physical Review Letters},
  volume  = {14},
  pages   = {57--59},
  year    = {1965}
}
@article{kerr1963,
  author  = {Kerr, Roy P.},
  title   = {Gravitational Field of a Spinning Mass},
  journal = {Physical Review Letters},
  volume  = {11},
  pages   = {237--238},
  year    = {1963}
}
@article{chandrasekhar1931,
  author  = {Chandrasekhar, S.},
  title   = {The Maximum Mass of Ideal White Dwarfs},
  journal = {The Astrophysical Journal},
  year    = {1931}
}
@article{oppenheimer1939,
  author  = {Oppenheimer, J. Robert and Snyder, Hartland},
  title   = {On Continued Gravitational Contraction},
  journal = {Physical Review},
  year    = {1939}
}
@article{friedmann1922,
  author  = {Friedmann, Alexander},
  title   = {\"Uber die Kr\"ummung des Raumes},
  journal = {Zeitschrift f\"ur Physik},
  year    = {1922}
}
\end{filecontents}
\addbibresource{sandbox-refs.bib}

\usepackage{citemsp}
\usepackage{amsthm}
\usepackage{array}
\usepackage{booktabs}
\usepackage[colorlinks=true,
            citecolor=linkblue,
            linkcolor=linkblue,
            urlcolor=linkblue,
            bookmarks=true]{hyperref}

\newtheorem{theorem}{Theorem}
\newtheorem{lemma}[theorem]{Lemma}

% Custom prefix registration
\citemspprefix{w}{w.}%   widget (arbitrary user prefix)

\begin{document}

\title{\texttt{citemsp} --- Unified Test Suite}
\author{Sandbox}
\date{\today}
\maketitle

\tableofcontents
\newpage

% ═══════════════════════════════════════════════════════════════════════════════
\section{Plain citations and range compression}
% ═══════════════════════════════════════════════════════════════════════════════

\noindent
\texttt{\textbackslash cite} baseline (5 consecutive):
\cite{einstein1915,schwarzschild1916,wald1984,hawking1974,misner1973} \\
\texttt{\textbackslash citemsp} same keys:
\citemsp{einstein1915, schwarzschild1916, wald1984, hawking1974, misner1973}

\medskip\noindent
Non-consecutive (\texttt{\textbackslash cite}):
\cite{einstein1915,schwarzschild1916,wald1984,misner1973,penrose1965,kerr1963} \\
Non-consecutive (\texttt{\textbackslash citemsp}):
\citemsp{einstein1915, schwarzschild1916, wald1984, misner1973, penrose1965, kerr1963}

\medskip\noindent
All ten:
\citemsp{einstein1915, schwarzschild1916, wald1984, hawking1974, misner1973,
         penrose1965, kerr1963, chandrasekhar1931, oppenheimer1939, friedmann1922} \\
Single: \citemsp{wald1984}


% ═══════════════════════════════════════════════════════════════════════════════
\section{Default locators: section (\textsection) and paragraph (\textparagraph)}
% ═══════════════════════════════════════════════════════════════════════════════

\noindent
Section only:              \citemsp{wald1984/6.1} \\
Section + paragraph:       \citemsp{wald1984/6.1/2} \\
Paragraph only (sub slot): \citemsp{wald1984//5} \\
Nested section number:     \citemsp{wald1984/3.2.1} \\
Deep subsection + para:    \citemsp{wald1984/3.2.1/4}


% ═══════════════════════════════════════════════════════════════════════════════
\section{Built-in prefix locators}
% ═══════════════════════════════════════════════════════════════════════════════

Each built-in prefix in both the superscript and subscript slots:

\medskip\noindent
\begin{tabular}{llll}
\textbf{Prefix} & \textbf{Type} & \textbf{In super slot} & \textbf{In sub slot} \\[2pt]
\texttt{d} (definition) & $\triangleq$  & \citemsp{wald1984/d5}        & \citemsp{wald1984/3.2/d5} \\
\texttt{e} (equation)   & eq.           & \citemsp{wald1984/e4.3}      & \citemsp{wald1984/3.2/e4.3} \\
\texttt{f} (figure)     & fig.          & \citemsp{wald1984/f3}        & \citemsp{wald1984/6.1/f3} \\
\texttt{n} (footnote)   & $\dagger$     & \citemsp{hawking1974/n7}     & \citemsp{hawking1974/1/n7} \\
\texttt{p} (page)       & pp.           & \citemsp{wald1984/p42}       & \citemsp{wald1984/3.2/p42} \\
\texttt{t} (table)      & $\boxplus$    & \citemsp{misner1973/t2}      & \citemsp{misner1973/8.4/t2} \\
\texttt{T} (theorem)    & Th.           & \citemsp{wald1984/T4.2}      & \citemsp{wald1984/3.2/T1} \\
\texttt{L} (lemma)      & Lem.          & \citemsp{wald1984/L1}        & \citemsp{wald1984/3.2/L1} \\
\texttt{C} (corollary)  & Cor.          & \citemsp{wald1984/C3}        & \citemsp{wald1984/3.2/C3} \\
\texttt{A} (appendix)   & App.          & \citemsp{wald1984/A2}        & \citemsp{wald1984/3.2/A2} \\
\end{tabular}

\medskip\noindent
Custom prefix (registered with \verb|\citemspprefix{w}{w.}|): \citemsp{misner1973/w7}

\medskip\noindent
Mixed built-in types in both slots: \citemsp{wald1984/T4.2/e3.1}


% ═══════════════════════════════════════════════════════════════════════════════
\section{Mixed locator types in a single call}
% ═══════════════════════════════════════════════════════════════════════════════

\noindent
\verb|\citemsp{wald1984/6.1/2, misner1973/p520, hawking1974/1/e4.3}|\\
Output: \citemsp{wald1984/6.1/2, misner1973/p520, hawking1974/1/e4.3}

\medskip\noindent
\verb|\citemsp{wald1984/T4.2, einstein1915/3.2/d5, misner1973/8.4/t2, kerr1963/p42}|\\
Output: \citemsp{wald1984/T4.2, einstein1915/3.2/d5, misner1973/8.4/t2, kerr1963/p42}


% ═══════════════════════════════════════════════════════════════════════════════
\section{Range compression with mixed entries}
% ═══════════════════════════════════════════════════════════════════════════════

\noindent
Range then locator:
\citemsp{einstein1915, schwarzschild1916, wald1984, hawking1974/3.2} \\
Locator then range:
\citemsp{einstein1915/3.2, schwarzschild1916, wald1984, hawking1974, misner1973} \\
Plain, locator, range:
\citemsp{einstein1915, schwarzschild1916/3.2, wald1984, hawking1974, misner1973} \\
Range, locator, range:
\citemsp{einstein1915, schwarzschild1916, wald1984, hawking1974/e2.1,
         misner1973, penrose1965, kerr1963} \\
Two locators between ranges:
\citemsp{einstein1915, schwarzschild1916, wald1984,
         hawking1974/3.1, misner1973/p23,
         penrose1965, kerr1963, chandrasekhar1931} \\
Adjacent locators only:
\citemsp{wald1984/3.1, hawking1974/e2.1, misner1973/p42}


% ═══════════════════════════════════════════════════════════════════════════════
\section{Italic, slanted, and bold-italic contexts}
% ═══════════════════════════════════════════════════════════════════════════════

\noindent
Upright: The field equations~\citemsp{wald1984/3.2/e4.3} can be derived from
a variational principle.

\medskip\noindent
\textit{Italic: Hawking radiation~\citemsp{hawking1974/1/2} implies black holes
are not truly black.}

\medskip\noindent
\textit{Italic with prefix: The ADM formalism~\citemsp{misner1973/p520}
decomposes spacetime into spatial hypersurfaces.}

\medskip\noindent
\textit{Italic with stacked locators: See~\citemsp{wald1984/6.1/e4.3} for the
derivation.}

\medskip\noindent
\textsl{Slanted: See the derivation~\citemsp{wald1984/3.2/e4.3} in the text.}

\medskip\noindent
\textbf{\textit{Bold italic: the metric tensor~\citemsp{wald1984/3.1/e3.4}
satisfies the symmetry.}}

\medskip\noindent
\textit{\textbf{Italic then bold: the metric tensor~\citemsp{wald1984/3.1/e3.4}
satisfies the symmetry.}}


% ═══════════════════════════════════════════════════════════════════════════════
\section{Font size contexts}
% ═══════════════════════════════════════════════════════════════════════════════

Each size command wraps \verb|\citemsp{wald1984/6.1/2}|:

\medskip\noindent
{\tiny         \texttt{tiny}:         \citemsp{wald1984/6.1/2} Lorem ipsum.}\\[2pt]
{\scriptsize   \texttt{scriptsize}:   \citemsp{wald1984/6.1/2} Lorem ipsum.}\\[2pt]
{\footnotesize \texttt{footnotesize}: \citemsp{wald1984/6.1/2} Lorem ipsum.}\\[2pt]
{\small        \texttt{small}:        \citemsp{wald1984/6.1/2} Lorem ipsum.}\\[2pt]
{\normalsize   \texttt{normalsize}:   \citemsp{wald1984/6.1/2} Lorem ipsum.}\\[2pt]
{\large        \texttt{large}:        \citemsp{wald1984/6.1/2} Lorem ipsum.}\\[2pt]
{\Large        \texttt{Large}:        \citemsp{wald1984/6.1/2} Lorem ipsum.}\\[2pt]
{\LARGE        \texttt{LARGE}:        \citemsp{wald1984/6.1/2} Lorem ipsum.}\\[2pt]
{\huge         \texttt{huge}:         \citemsp{wald1984/6.1/2} Lorem ipsum.}


% ═══════════════════════════════════════════════════════════════════════════════
\section{Scaling: \texttt{\textbackslash citemspscale}}
% ═══════════════════════════════════════════════════════════════════════════════

\noindent
Default scale (0.35): \citemsp{wald1984/6.1/e4.3}

{%
\renewcommand{\citemspscale}{0.25}%
Smaller (0.25): \citemsp{wald1984/6.1/e4.3}
}

{%
\renewcommand{\citemspscale}{0.50}%
Larger (0.50): \citemsp{wald1984/6.1/e4.3}
}


% ═══════════════════════════════════════════════════════════════════════════════
\section{Small caps and \texttt{\textbackslash emph} toggle}
% ═══════════════════════════════════════════════════════════════════════════════

\noindent
\textsc{Small caps: The Birkhoff theorem~\citemsp{wald1984/T6.2} is a
cornerstone of GR.}

\medskip\noindent
Upright: \emph{The metric~\citemsp{wald1984/3.1} is a fundamental object.}

\medskip\noindent
\textit{Italic outer: \emph{upright inner~\citemsp{wald1984/3.2/e3.4}} with
toggled shape.}

\medskip\noindent
\textit{Italic: The curvature~\citemsp{wald1984/3.2/e3.4} satisfies the
Bianchi identity. Inside it: \emph{the Riemann
tensor~\citemsp{wald1984/3.2}} is antisymmetric.}


% ═══════════════════════════════════════════════════════════════════════════════
\section{Math mode adjacency}
% ═══════════════════════════════════════════════════════════════════════════════

The citation should survive adjacent math without mode-switch errors.

\medskip\noindent
$E = mc^2$~\citemsp{einstein1915/3.2}

\medskip\noindent
\citemsp{einstein1915/e1}~implies $G_{\mu\nu} = 8\pi T_{\mu\nu}$.

\medskip\noindent
Between two math spans:
$\nabla_\mu T^{\mu\nu}=0$~\citemsp{wald1984/4.3/e4.3}~gives
$\dot\rho + 3H(\rho+p)=0$.


% ═══════════════════════════════════════════════════════════════════════════════
\section{Theorem and list environments}
% ═══════════════════════════════════════════════════════════════════════════════

\begin{theorem}[Birkhoff \citemsp{wald1984/T6.2}]
  Every spherically symmetric solution of the vacuum Einstein field equations
  is isometric to the Schwarzschild metric~\citemsp{wald1984/6.1/e6.1}.
\end{theorem}

\begin{lemma}
  The area of the event horizon is
  non-decreasing~\citemsp{hawking1974/1/e2}.
\end{lemma}

\subsection*{Itemize}

\begin{itemize}
  \item The field equations~\citemsp{einstein1915} were derived in 1915.
  \item The Schwarzschild solution~\citemsp{wald1984/6.1/2} followed shortly.
  \item Gravitational waves are a prediction of linearised
        GR~\citemsp{misner1973/p35.1}.
\end{itemize}

\subsection*{Description}

\begin{description}
  \item[Singularity theorems] proved by
    Penrose~\citemsp{penrose1965/T1} and others.
  \item[Hawking radiation] derived
    in~\citemsp{hawking1974/e1/n2}.
  \item[Kerr metric]
    see~\citemsp{kerr1963/3.2}.
\end{description}


% ═══════════════════════════════════════════════════════════════════════════════
\section{Table cells}
% ═══════════════════════════════════════════════════════════════════════════════

\medskip\noindent
\begin{tabular}{lll}
  \toprule
  \textbf{Source}         & \textbf{Year} & \textbf{Citation} \\
  \midrule
  Field equations         & 1915          & \citemsp{einstein1915/3.2} \\
  General Relativity book & 1984          & \citemsp{wald1984/6.1/2} \\
  Gravitational collapse  & 1965          & \citemsp{penrose1965/T1} \\
  Hawking radiation       & 1974          & \citemsp{hawking1974/e1} \\
  \bottomrule
\end{tabular}


% ═══════════════════════════════════════════════════════════════════════════════
\section{Long locator strings and many-key calls}
% ═══════════════════════════════════════════════════════════════════════════════

Deep hierarchical section number: \citemsp{wald1984/12.3.4.5}

\medskip\noindent
Both slots with a long equation number: \citemsp{wald1984/12.3.4.5/e99.999}

\medskip\noindent
Long page reference: \citemsp{misner1973/p1279}

\medskip\noindent
Six keys mixing plain and locators:
\citemsp{wald1984/6.1/2, einstein1915, hawking1974, misner1973/p520,
         penrose1965, kerr1963/3.2}

\medskip\noindent
Six keys, all plain (maximum compression):
\citemsp{wald1984, einstein1915, hawking1974, misner1973, penrose1965, kerr1963}


% ═══════════════════════════════════════════════════════════════════════════════
\section{Double-slash / extra-slash edge cases}
% ═══════════════════════════════════════════════════════════════════════════════

The parser is \verb|\@citemsp@parse#1/#2/#3/#4\@nil|; extra slashes may land
in \verb|#4| (discarded) or produce empty tokens in \verb|#2|/\verb|#3|.

\medskip\noindent
\textbf{Empty second field (subscript-only):}
\verb|\citemsp{wald1984//5}| $\to$ \citemsp{wald1984//5}

\medskip\noindent
\textbf{Three slashes (key + empty sec + empty par + trailing):}
\verb|\citemsp{wald1984///}| $\to$ \citemsp{wald1984///}

\medskip\noindent
\textbf{Four-part with trailing slash (extra slash in \#4):}
\verb|\citemsp{wald1984/6.1/2/extra}| $\to$ \citemsp{wald1984/6.1/2/extra}

\medskip\noindent
\textbf{Only slashes, no key (degenerate):}
\verb|\citemsp{///}| $\to$ \citemsp{///}


% ═══════════════════════════════════════════════════════════════════════════════
\section{Footnotes}
% ═══════════════════════════════════════════════════════════════════════════════

Here is a running-text reference.\footnote{See~\citemsp{wald1984/3.2} for
the full derivation, and also~\citemsp{hawking1974/1} for the original
energy-condition argument.}


% ═══════════════════════════════════════════════════════════════════════════════
\section{Drop-in \texttt{\textbackslash cite} replacement and realistic text}
% ═══════════════════════════════════════════════════════════════════════════════

\texttt{\textbackslash citemsp} with no slashes behaves like
\texttt{\textbackslash cite}:

\medskip\noindent
\verb|\cite{wald1984}| $\longrightarrow$ \cite{wald1984} \\
\verb|\citemsp{wald1984}| $\longrightarrow$ \citemsp{wald1984} \\[4pt]
\verb|\cite{einstein1915,...,misner1973}| $\longrightarrow$
\cite{einstein1915,schwarzschild1916,wald1984,hawking1974,misner1973} \\
\verb|\citemsp{einstein1915,...,misner1973}| $\longrightarrow$
\citemsp{einstein1915, schwarzschild1916, wald1984, hawking1974, misner1973}

\bigskip\noindent
\textbf{Realistic paragraph:}
The Schwarzschild solution~\citemsp{schwarzschild1916/1/3} was found shortly
after the field equations~\citemsp{einstein1915} were published. A modern
treatment can be found in standard
textbooks~\citemsp{wald1984/6.1/2, misner1973/p520}, where the derivation
proceeds via the Birkhoff theorem~\citemsp{wald1984/T6.2}. The singularity
theorems~\citemsp{penrose1965, hawking1974/1/e4.3} later showed that
gravitational collapse is generic. For further context
see~\citemsp{kerr1963, chandrasekhar1931, oppenheimer1939, friedmann1922}.


% ═══════════════════════════════════════════════════════════════════════════════
\printbibliography
% ═══════════════════════════════════════════════════════════════════════════════

\end{document}
